
Haben wir $\C^*$-Algebren $\Aa, \Ba$ mit respektiven Abbildungen $\varepsilon_\Aa, \varepsilon_\Ba$, so können wir schreiben
\[
    \varepsilon_{\Ba}^{-1}(b) = \sum_{n=1}^\infty b^n.
\]
Damit ist also
\[
    \restr{\varepsilon_\Ba^{-1}}{\Aa_+ \cap B_1^\Aa (0)} = \varepsilon_\Aa^{-1},
\]
womit folgt
\[
    \varepsilon_\Aa = \restr{\varepsilon_\Ba}{\Aa_+ \cap B_1^\Aa (0)}.
\]

\begin{lemma}
    Sei $\Aa$ eine $C^*$-Algebra, so ist das Netz $(\varepsilon(q))_{q \in \Aa_+}$ (mit Werten in $\Aa_+ \cap B_1^\Aa (0)$) eine approximative Eins, das heißt
    \[
        \lim_{q \in \Aa_+} \varepsilon(q) a = \lim_{q \in \Aa_+} a \varepsilon(q) = a, \quad \forall a \in \Aa.
    \]
\end{lemma}

\begin{proof}
    Wir erinnern, dass wir für Elemente $p, q \in \Aa_+$ bereits gezeigt haben:
    \begin{itemize}
        \item $p \leq q \implies \Vert p \Vert \leq \Vert q \Vert$
        \item $p \leq \Vert p \Vert e$
        \item $p \leq e \Leftrightarrow \Vert p \Vert \leq 1$
        \item Für $c \in \Aa$ und $p \leq q$ gilt $c^* p c \leq c^* q c$.
    \end{itemize}
    Es reicht nun die Aussage für $a \in \Aa_+$ zu zeigen, da wir ein allgemeines $a \in \Aa$ in Real- und Imaginärteil zerlegen können (welche selbstadjungiert sind), und anschließend die selbstadjungierten Elemente in Differenz zweier positiver Elemente. Vermöge einer Streckung können wir sogar $\Vert a \Vert < 1$ annehmen.

    Sei nun $\eta > 0$ und setze $q_\eta \coloneqq \frac{1}{\eta^2} a \in \Aa_+$. Für beliebiges $q \in \Aa_+, q \geq q_\eta$ gilt nun $\varepsilon(q_\eta) \leq \varepsilon(q) \leq e$ und $e - \varepsilon(q) \leq e$. Damit folgt
    \begin{align*}
        0
        &\leq
        (e - \varepsilon(q))^2
        =
        \sqrt{e - \varepsilon(q)} (e - \varepsilon(q)) \sqrt{e - \varepsilon(q)}
        \leq
        e - \varepsilon(q) \\
        &\leq
        e - \varepsilon(q_\eta)
        = \left( e + \frac{1}{\eta^2} a \right)^{-1}.
    \end{align*}
    Konjugieren mit $a$ liefert
    \begin{align*}
        0
        &\leq
        a(e - \varepsilon(q))^2 a
        \leq
        \sqrt{a} \left[ \sqrt{a} \left( e + \frac{1}{\eta^2} a \right)^{-1} \sqrt{a} \right] \sqrt{a}.
    \end{align*}
    Der Term in der $[\cdots]$ Klammer ist gleich
    \[
        a \left( e + \frac{1}{\eta^2} a \right)^{-1} = \eta^2 \frac{1}{\eta^2} a \left( e + \frac{1}{\eta^2} a \right)^{-1} = \eta^2 \varepsilon(q_\eta) \leq \eta^2 e,
    \]
    womit insgesamt die Ungleichung
    \[
        0
        \leq
        a (e - \varepsilon(q))^2 a \leq \eta^2 a
    \]
    folgt. Anwenden der Norm liefert
    \[
        \Vert a (e - \varepsilon(q))^2 a \Vert \leq \eta^2 \Vert a \Vert < \eta^2.
    \]
    Damit folgt nun
    \begin{align*}
        \Vert a - \varepsilon(q) a \Vert^2
        &=
        \Vert a - a \varepsilon(q) \Vert^2
        =
        \Vert a (e - \varepsilon(q) ) \Vert^2
        =
        \Vert a (e - \varepsilon(q))^2 a \Vert
        <
        \eta^2.
    \end{align*}
    Damit folgt die erste Gleichheit in der Aussage, und aus Symmetriegründen auch der Rest.
\end{proof}

\begin{proposition}
    Sei $\Aa$ eine $C^*$-Algebra und $J$ ein echtes, abgeschlossenes Ideal von $\Aa$. Dann gilt:
    \begin{enumerate}
        \item $J^* = J$. Insbesondere ist $J$ eine $C^*$-Unteralgebra.
        \item Die Faktorraumnorm auf $A/_J$ erfüllt
        $$ \Vert [a] \Vert = \lim_{q \in J_+} \Vert a - a \varepsilon(q) \Vert. $$
        \item $A/_J$ versehen mit der Faktorraumnorm ist eine $C^*$-Algebra, falls $[a]^* \coloneqq [a^*]$.
    \end{enumerate}
\end{proposition}

\begin{proof}{\ }
    \begin{enumerate}
        \item Sei $d \in J$ und betrachte die (in $\Aa$) erzeugte $C^*$-Algebra
        $$ \Ca \coloneqq [{d^* d}] = \clspn \{ (d^* d)^n : n \in \N \} \leq \Aa. $$
        Für $q \in \Ca_+$ gilt $\varepsilon(q) \in \Ca_+$ und
        \begin{align*}
            \Vert \varepsilon(q) d^* - d^* \Vert^2
            &=
            \Vert d \varepsilon(q) - d \Vert^2
            =
            \Vert d (\varepsilon(q) - e) \Vert^2
            =
            \Vert (\varepsilon(q) - e) d^* d (\varepsilon(q) - e) \Vert \\
            &=
            \Vert \varepsilon(q) - e \Vert \, \Vert d^* d \varepsilon(q) - d^* d \Vert
            \leq
            2 \Vert d^* d \varepsilon(q) - d^* d \Vert.
        \end{align*}
        Mit der approximativen Einheit erhalten wir als Grenzwert für $q \in \Ca_+$, dass der obige Ausdruck gegen $0$ konvergiert. Also ist $d^* \in \cl{J} = J$.

        \item Sei $a \in \Aa, d \in J, q \in J_+$, dann ist $a \varepsilon(q) \in J$ und $0 \leq e - \varepsilon(q) \leq e$. Jetzt gilt
        \begin{align*}
            \Vert [a] \Vert
            &=
            \inf_{j \in J} \Vert a - j \Vert
            =
            \Vert [a - a \varepsilon(q) ] \Vert
            \leq 
            \Vert a - a \varepsilon(q) \Vert \\
            &\leq
            \Vert (a - d)(e - \varepsilon(q)) \Vert + \Vert d - d \varepsilon(q) \Vert
            \leq
            \Vert a - d \Vert + \Vert d - d \varepsilon(q) \Vert.
        \end{align*}
        Nehmen wir den Limes superior, so erhalten wir
        \begin{align*}
            \Vert [a] \Vert
            &\leq
            \limsup_{q \in J_+} \Vert a - a \varepsilon(q) \Vert
            \leq
            \Vert a - d \Vert, \quad \forall a \in \Aa, d \in J.
        \end{align*}
        Anwenden des Infimums über $d \in J$ liefert
        \[
            \Vert [a] \Vert
            \leq
            \limsup_{q \in J_+} \Vert a - a \varepsilon(q) \Vert
            \leq
            \Vert [a] \Vert
            \leq
            \liminf_{q \in J_+} \Vert a - a \varepsilon(q) \Vert
        \]
        Insbesondere existiert der Limes und stimmt mit der Faktorraumnorm überein.

        \item Zunächst gilt
        \begin{align*}
            \Vert (e - \varepsilon(q)) a^* a (e - \varepsilon(q)) \Vert = \Vert a - a \varepsilon(q) \Vert^2,
        \end{align*}
        wobei wegen dem vorigen Punkt die Limiten über $q \in J_+$ existieren und mit $\Vert [a] \Vert^2$ übereinstimmen. Damit gilt
        \begin{align*}
            \Vert [a]^* [a] \Vert
            &=
            \Vert [a^* a] \Vert
            =
            \lim_{q \in J_+} \Vert a^* a - a^* a \varepsilon(q) \Vert
            =
            \lim_{q \in J_+} \Vert a^* a (e - \varepsilon(q)) \Vert \\
            &\geq
            \lim_{q \in J_+} \Vert (e - \varepsilon(q)) a^* a (e - \varepsilon(q)) \Vert
            =
            \Vert [a] \Vert^2. \qedhere
        \end{align*}
    \end{enumerate}
\end{proof}

\begin{example}
    Sei $H$ ein unendlichdimensionaler Hilbertraum und betrachte die kompakten Operatoren $K(H)$ als echtes Unterideal von $\Aa \coloneqq L_b(H)$. Nach dem oben Gezeigten ist $L_b(H)/_{K(H)}$ eine $C^*$-Algebra.
\end{example}

\begin{corollary}
    Seien $\Aa, \Ba$ $C^*$-Algebren und sei $\varphi$ ein $\ast$-Algebrenhomomorphismus von $\Aa$ nach $\Ba$, $\varphi \neq 0$. Dann bildet $\varphi(\Aa)$ eine Unter-$C^*$-Algebra von $\Ba$ und
    \[
        \varphi/_{\ker \varphi} : \Aa/_{\ker \varphi} \to \varphi(\Aa)
    \]
    ist ein isometrischer $\ast$-Algebrenisomorphismus.
\end{corollary}

\begin{proof}
    Zunächst ist $\varphi$ stetig (sogar eine Kontraktion), wie wir bereits gezeigt haben. Weiters ist $\ker \varphi$ ein echtes, abgeschlossenes Ideal in $\Aa$. Nach der vorigen Proposition ist $\Aa_{\ker \varphi}$ eine $C^*$-Algebra. Die Abbildung
    \[
        \mapping{\varphi/_{\ker \varphi}}
            {\Aa/_{\ker \varphi} & \to & \varphi(\Aa),}
            {\textrm{[}a\textrm{]} & \mapsto & \varphi(a)}
    \]
    ist dann ein injektiver $\ast$-Algebrenhomomorphismus. Da $\varphi/_{\ker \varphi}$ auch isometrisch ist, ist das Bild abgeschlossen, womit $\varphi(\Aa)$ eine $C^*$-Algebra ist.
\end{proof}

Wir erinnern, dass ein Funktional $f \in \Aa^*$ \emph{positiv} hieß, wenn $f(q) \geq 0$ für alle $q \in \Aa_+$, wobei $\Aa$ eine $C^*$-Algebra mit Eins war. Wir wissen bereits, dass positive $f$ beschränkt sind, also $\Vert f \Vert < \infty$. Außerdem ist $f$ positiv genau dann wenn $f(e) = \Vert f \Vert$.

Wir weitern dieses Konzept aus:

\begin{definition}
    Sei $\Aa$ eine $C^*$-Algebra, dann heißt ein $f \in \Aa^*$ positiv, falls
    $$ f(q) \geq 0, \quad \forall q \in \Aa_+. $$
\end{definition}

\begin{definition}
    Sei $f \in \Aa^*$ positiv. Dann definieren wir
    \[
        f(e_-) \coloneqq \lim_{q \in \Aa_+} f(\varepsilon(q)) = \sup_{q \in \Aa_+} f(\varepsilon(q)) = \sup_{a \in \Aa_+ \cap B_1^\Aa (0)} f(a) \leq \Vert f \Vert.
    \]
\end{definition}
